\documentclass[a4paper,11pt]{article}
\usepackage[a4paper,margin=2cm]{geometry}
\usepackage[T1]{fontenc}
\usepackage[utf8]{inputenc}
\usepackage[ngerman,english]{babel}
\usepackage{lmodern}
\usepackage[hidelinks]{hyperref}
\title{Freier Dativ: Dativus commodi und incommodi}\date{}\begin{document}\maketitle\tableofcontents
\section{Was ist ein freier Dativ?}
Ein freier Dativ ist eine zusätzliche Dativ-Ergänzung, die nicht in gleicher Weise vom Verb gefordert wird wie ein obligatorisches Dativobjekt. Der Dativ bezeichnet z.\,B. die Person, zu deren Vorteil oder Nachteil etwas geschieht.
\textbf{English:} A free dative is an optional dative expression indicating, among other things, who benefits or suffers from an action.

\section{Dativus commodi (Dativ des Vorteils)}
Er nennt die Person, \textbf{zu deren Vorteil} eine Handlung geschieht.
\begin{itemize}
\item \textit{Ich kaufe mir ein Fahrrad.} -- I buy myself a bicycle.
\item \textit{Du mietest dir ein Auto.} -- You rent a car for yourself.
\item \textit{Sie kocht ihrem Kind eine Suppe.} -- She cooks soup for her child.
\end{itemize}
\textbf{Struktur:} Subjekt + Verb + Dativ (Begünstigter) + Akkusativobjekt (Sache), sofern das Verb ein Akkusativobjekt hat. Der Dativ ist nicht immer notwendig: \textit{Ich miete ein Auto} ist ebenfalls richtig.

\section{Dativus incommodi (Dativ des Nachteils)}
Er bezeichnet die Person, \textbf{zu deren Nachteil} etwas geschieht.
\begin{itemize}
\item \textit{Jemand hat mir das Fahrrad gestohlen.} -- Someone stole my bicycle (to my detriment).
\item \textit{Du hast mir die Tasse zerbrochen.} -- You broke my cup (affecting me).
\end{itemize}
Der genaue Status kann je nach Verb und Analyse unterschiedlich beschrieben werden; der Dativ drückt hier Betroffenheit aus.

\section{Reflexiver Dativ: mir, dir, sich, uns, euch, sich}
Wenn Dativ und Subjekt dieselbe Person bezeichnen, werden die reflexiven Formen verwendet:
\begin{center}
\begin{tabular}{ll}
ich & mir\\ du & dir\\ er/sie/es & sich\\ wir & uns\\ ihr & euch\\ sie/Sie & sich
\end{tabular}
\end{center}
\textit{Ich miete mir ein Zimmer.} -- \textit{Wir mieten uns ein Zimmer.}

\textbf{Wichtig:} \textit{mieten} ist \emph{nicht} obligatorisch reflexiv. Hier steht ein freier/reflexiver Dativ; man kann ihn weglassen: \textit{Ich miete ein Zimmer.}

\section{Vergleich: für + Akkusativ}
\textit{Ich kaufe mir eine Jacke.} $\approx$ \textit{Ich kaufe eine Jacke für mich.}

\textit{für mich} betont die Zweckbestimmung bzw. den Begünstigten stärker; nicht in jedem Kontext ist es gleichwertig mit einem freien Dativ.
\textit{Du mietest dir das Auto.} $\approx$ \textit{Du mietest das Auto für dich.}
\textbf{Beachte:} \textit{das Auto} und \textit{ein Auto} unterscheiden sich durch Bestimmtheit.

\section{Abgrenzung vom obligatorischen Dativobjekt}
\textit{Ich helfe dir.} -- \textit{helfen} verlangt standardsprachlich einen Dativ; \textit{dir} ist hier kein freier Dativus commodi.\\
\textit{Ich gebe dir das Buch.} -- \textit{dir} ist die Empfängerergänzung von \textit{geben}; nicht einfach ein frei hinzugefügter Dativ des Vorteils.\\
\textit{Ich kaufe mir ein Buch.} -- \textit{mir} ist ein freier Dativ des Vorteils.

\section{Weitere freie Dative (Überblick)}
\textbf{Dativus possessivus / Pertinenzdativ:} \textit{Ich wasche mir die Hände.} (Zuordnung zum Körper / Besitzer; nicht automatisch Dativus commodi.)\\
\textbf{Dativus ethicus:} \textit{Pass mir bloß auf!} (sprecherbezogene Anteilnahme; umgangssprachlich).
Diese Typen unterscheiden sich in Bedeutung und Grammatik.

\section{Übung}
Bestimme die Funktion des Dativs: (1) \textit{Ich backe mir Brot.} (2) \textit{Er hilft mir.} (3) \textit{Ich putze mir die Zähne.} (4) \textit{Jemand hat mir den Schlüssel weggenommen.}

\textbf{Lösungen:} (1) commodi; (2) vom Verb verlangtes Dativobjekt; (3) Pertinenzdativ; (4) Dativ der betroffenen/benachteiligten Person (incommodi).
\end{document}
